\chapter{TRIỂN KHAI, KIỂM THỬ VÀ ĐÁNH GIÁ HỆ THỐNG}

```latex
\section{Môi trường triển khai}

Để triển khai Notification Service, nhóm B7 sử dụng môi trường phát triển gồm Windows 11 kết hợp với Docker Desktop nhằm quản lý và vận hành các container của hệ thống. Việc sử dụng Docker giúp đảm bảo tính đồng nhất giữa môi trường phát triển và môi trường triển khai thực tế, đồng thời giảm thiểu các vấn đề liên quan đến cấu hình phần mềm.

Các công nghệ và công cụ được sử dụng trong quá trình phát triển được trình bày trong Bảng \ref{tab:environment}.

\begin{table}[H]
\centering
\caption{Môi trường phát triển}
\label{tab:environment}
\begin{tabular}{|l|l|}
\hline
Hệ điều hành & Windows 11 \\ \hline
Ngôn ngữ lập trình & Python 3.11 \\ \hline
Framework & FastAPI \\ \hline
Cơ sở dữ liệu & PostgreSQL 15 \\ \hline
Container & Docker Desktop \\ \hline
Quản lý Container & Docker Compose \\ \hline
Công cụ kiểm thử & PowerShell, Postman \\ \hline
Kênh gửi thông báo & Telegram Bot API \\ \hline
Quản lý mã nguồn & GitHub \\ \hline
\end{tabular}
\end{table}

Trong quá trình triển khai, Notification Service được cấu hình để giao tiếp với PostgreSQL và Mock Channel Service thông qua mạng nội bộ Docker. Đồng thời, Notification Service cũng kết nối với Telegram Bot API để gửi cảnh báo tới người dùng.

\section{Triển khai bằng Docker Compose}

Notification Service được triển khai bằng Docker Compose nhằm quản lý đồng thời nhiều container.

Sau khi thực hiện lệnh:

\begin{center}
\texttt{docker compose up -d --build}
\end{center}

hệ thống tự động khởi tạo ba container gồm:

\begin{itemize}
\item Notification API
\item PostgreSQL Database
\item Mock Channel Service
\end{itemize}

Sau khi khởi động thành công, trạng thái các container được kiểm tra bằng lệnh:

\begin{center}
\texttt{docker compose ps}
\end{center}

Kết quả cho thấy toàn bộ container đều chuyển sang trạng thái Healthy, chứng tỏ Notification Service đã được triển khai thành công.

\begin{figure}[H]
\centering
\fbox{\parbox{0.9\textwidth}{\centering Chèn ảnh docker compose ps}}
\caption{Kết quả triển khai Docker Compose}
\label{fig:docker-ps}
\end{figure}

\section{Kiểm thử Health Check}

Sau khi triển khai thành công, nhóm tiến hành kiểm tra trạng thái hoạt động của Notification Service thông qua endpoint:

\begin{center}
\texttt{GET /health}
\end{center}

API trả về thông tin về trạng thái của Notification Service, PostgreSQL và Mock Channel Service.

Kết quả kiểm thử cho thấy cả ba thành phần đều hoạt động bình thường và sẵn sàng phục vụ các yêu cầu từ Core Business Service.

\begin{figure}[H]
\centering
\fbox{\parbox{0.9\textwidth}{\centering Chèn ảnh kết quả GET /health}}
\caption{Kết quả Health Check}
\label{fig:health}
\end{figure}

\section{Kiểm thử chức năng gửi thông báo}

Để kiểm tra chức năng chính của Notification Service, nhóm sử dụng PowerShell gửi yêu cầu POST đến endpoint:

\begin{center}
\texttt{/api/notifications/alerts}
\end{center}

Request bao gồm các thông tin:

\begin{itemize}
\item Alert ID
\item Severity
\item Message
\item Target
\item Channel
\end{itemize}

Sau khi tiếp nhận yêu cầu, Notification Service thực hiện:

\begin{enumerate}
\item Xác thực Bearer Token.
\item Validate dữ liệu đầu vào.
\item Kiểm tra Alert ID.
\item Gửi Telegram.
\item Lưu lịch sử vào PostgreSQL.
\item Trả Response về cho Client.
\end{enumerate}

Kết quả kiểm thử cho thấy Notification Service trả về trạng thái:

\begin{center}
Delivered
\end{center}

và Telegram Bot nhận được thông báo thành công.

\begin{figure}[H]
\centering
\fbox{\parbox{0.9\textwidth}{\centering Chèn ảnh PowerShell POST}}
\caption{Kiểm thử API gửi thông báo}
\label{fig:post-alert}
\end{figure}

\begin{figure}[H]
\centering
\fbox{\parbox{0.9\textwidth}{\centering Chèn ảnh Telegram}}
\caption{Thông báo nhận được trên Telegram}
\label{fig:telegram}
\end{figure}

\section{Kiểm thử API Logs}

Sau khi gửi thành công thông báo, nhóm tiến hành kiểm tra lịch sử xử lý thông qua endpoint:

\begin{center}
\texttt{/internal/notifications/logs}
\end{center}

API trả về danh sách các thông báo đã xử lý bao gồm:

\begin{itemize}
\item Alert ID
\item Severity
\item Message
\item Status
\item Attempts
\item Channel
\item Detail
\item Created Time
\end{itemize}

Kết quả cho thấy dữ liệu được lưu chính xác trong PostgreSQL và có thể truy xuất thông qua REST API.

\begin{figure}[H]
\centering
\fbox{\parbox{0.9\textwidth}{\centering Chèn ảnh Logs API}}
\caption{Kết quả truy xuất lịch sử thông báo}
\label{fig:logs}
\end{figure}

\section{Kiểm thử tích hợp giữa các nhóm}

Một trong những yêu cầu quan trọng của bài tập lớn là khả năng tích hợp giữa các nhóm phát triển.

Notification Service của nhóm B7 đã được tích hợp thành công với:

\subsection{Tích hợp với Core Business Service (B6)}

Core Business Service gửi Alert đến Notification Service thông qua REST API.

Notification Service tiếp nhận yêu cầu, xử lý và trả kết quả thành công.

Quá trình kiểm thử được thực hiện trên mạng LAN với địa chỉ IP của nhóm B7.

\subsection{Tích hợp với Analytics Service (B5)}

Analytics Service sử dụng API:

\begin{center}
\texttt{/internal/notifications/logs}
\end{center}

để truy xuất lịch sử gửi thông báo.

Các dữ liệu này được sử dụng để xây dựng Dashboard thống kê và đánh giá hiệu quả hoạt động của hệ thống.

\section{Đánh giá kết quả}

Sau quá trình triển khai và kiểm thử, Notification Service đã hoàn thành đầy đủ các chức năng theo yêu cầu của bài tập lớn.

Các chức năng đã hoàn thành bao gồm:

\begin{itemize}
\item Triển khai bằng Docker Compose.
\item Health Check.
\item Bearer Authentication.
\item Validate Request.
\item Retry Mechanism.
\item Idempotency.
\item Telegram Integration.
\item PostgreSQL Logging.
\item RESTful API.
\item OpenAPI Contract.
\item Tích hợp thành công với nhóm B6 và nhóm B5.
\end{itemize}

Kết quả kiểm thử thực tế cho thấy Notification Service hoạt động ổn định, thời gian phản hồi nhanh và đáp ứng đầy đủ yêu cầu của hệ thống Smart Campus.

\section{Hạn chế và hướng phát triển}

Mặc dù đã hoàn thành các yêu cầu của bài tập lớn, Notification Service vẫn còn một số hạn chế.

Hiện tại hệ thống mới chỉ triển khai kênh Telegram. Các kênh gửi khác như Email, SMS hoặc Discord vẫn chưa được phát triển.

Ngoài ra, hệ thống chưa xây dựng Dashboard quản trị trực quan để theo dõi trạng thái gửi thông báo theo thời gian thực.

Trong tương lai, Notification Service có thể được mở rộng theo các hướng sau:

\begin{itemize}
\item Hỗ trợ nhiều kênh gửi thông báo.
\item Xây dựng Dashboard quản trị.
\item Triển khai Message Queue (RabbitMQ hoặc Kafka).
\item Hỗ trợ WebSocket để cập nhật trạng thái thời gian thực.
\item Bổ sung cơ chế phân quyền người dùng.
\item Tăng cường khả năng giám sát bằng Prometheus và Grafana.
\end{itemize}
```
